\documentclass[11pt,a4paper]{article}

\usepackage[T1]{fontenc}
\usepackage[utf8]{inputenc}
\usepackage{lmodern}
\usepackage[margin=2.5cm]{geometry}
\usepackage{amsmath,amssymb}
\usepackage{booktabs}
\usepackage{array}
\usepackage{microtype}
\usepackage[authoryear,round]{natbib}
\usepackage{url}
\usepackage[hidelinks]{hyperref}

\title{Self-Supervised Network Reconstruction of Daily GNSS Time Series in California:\\
A Benchmark for Signal-Preserving Common-Mode Removal and Transient Detection}

\author{H. O\u{g}uz Bolat\\[2pt]
\small Department of Geomatics Engineering, Y{\i}ld{\i}z Technical University, Istanbul, T\"urkiye\\
\small \texttt{oguz.bolat@std.yildiz.edu.tr}}

\date{Research proposal, October 2026}

\begin{document}

\maketitle

\begin{abstract}
Daily Global Navigation Satellite System (GNSS) position time series resolve tectonic deformation at the millimetre level, yet their utility for detecting aseismic transients is limited by spatially correlated noise that is comparable in amplitude to the signals of interest. Classical common-mode filters and recent deep-learning denoisers reduce this noise, but they have been evaluated predominantly by the reduction of residual scatter, a metric that rewards the removal of genuine signal. Where the recovery of synthetic or independently observed signals has been assessed, it has not been examined as a function of spatial scale or compared with linear filters under a common protocol, and the multi-station deep-learning denoisers have been developed almost exclusively for subduction zones. We propose a self-supervised reconstruction task for a dense network of 1,381 continuously operating stations in and around California processed by the Nevada Geodetic Laboratory: each station's displacement is withheld and predicted from its neighbours, so that the prediction captures the spatially coherent part of the signal and the residual isolates station-specific content. Because real transients are themselves spatially coherent, neighbours within an exclusion radius are withheld as well, and this radius is varied systematically to quantify the trade-off between noise suppression and signal leakage. A hierarchy of models of increasing capacity, ranging from distance-weighted stacking and robust neighbour medians to per-station ridge regression, frequency-dependent filtering and a graph-attention neural network, will be trained and evaluated under identical temporal and spatial hold-out protocols. Performance will be assessed not only by reconstruction error but by the recovery of synthetic transients injected into real noise and by the preservation of the 2019 Ridgecrest postseismic signal. The study is designed to determine whether, and under which conditions, deep networks improve upon linear spatial filtering, and to provide an open benchmark and a calibrated noise model on which continual, near-real-time transient detection can subsequently be built.
\end{abstract}

\section{Introduction}

Continuous GNSS networks record the position of thousands of geodetic monuments at daily resolution with formal uncertainties of approximately 1~mm in the horizontal and 3--5~mm in the vertical. The Nevada Geodetic Laboratory (NGL) processes more than 17,000 such stations worldwide with a uniform strategy and distributes the resulting time series openly \citep{blewitt2018,ngl2026}, which has made network-scale analyses of crustal deformation routinely possible. Beyond steady interseismic motion, these records contain coseismic offsets, postseismic relaxation, slow slip events, magmatic inflation and hydrological loading, as well as instrumental artefacts such as antenna and receiver changes. Separating these contributions is the central problem of geodetic time-series analysis.

The dominant obstacle is noise that is correlated across stations. Errors in satellite orbits and clocks, reference-frame realisation and large-scale environmental loading produce a common-mode component that can reach several millimetres and that is coherent over hundreds to thousands of kilometres \citep{wdowinski1997,dong2006,kreemer2021}. Because this component is spatially coherent, it is readily confused with tectonic transients of comparable wavelength. The controversy surrounding a reported precursory acceleration in stacked high-rate GNSS data before large earthquakes \citep{bletery2023} illustrates the consequences: the signal disappears once a common-mode correction is applied, and the disagreement that followed concerns precisely whether such corrections remove noise only or also remove genuine deformation \citep{bradley2023,bletery2025}.

The present proposal arises from a set of preliminary analyses carried out with NGL daily solutions. Departures from a trajectory model fitted to a reference interval were examined for stations surrounding the East Anatolian Fault between the 2020 Elaz{\i}\u{g} and 2023 Kahramanmara\c{s} earthquakes and for eleven continuous stations within 80~km of the 2019 Ridgecrest sequence. In T\"urkiye the publicly available record is fragmentary, with data gaps spanning 2016--2020 and 2021--2022, which precludes a continuous interseismic analysis. At Ridgecrest, where the record is continuous, the coseismic offsets reach 0.63~m in the near field and postseismic displacement of 10--25~mm accumulates over the following two years. In the months preceding the mainshock, however, horizontal departures remain below 1~mm in the east component, while the north component shifts by 0.3--1.7~mm at every station with the same sign, irrespective of distance. A coherent vertical depression of 3--6~mm during early 2019 coincides with an exceptionally wet winter and recovers subsequently. Every apparent anomaly identified in these analyses is therefore attributable to network-wide noise, hydrological loading or equipment changes rather than to fault-local deformation. These observations motivate the present study, which shifts the emphasis from the search for precursors to the prior problem of separating coherent noise from coherent signal.

A natural formulation of this separation is to predict each station from its neighbours. The prediction can reproduce only those components that the station shares with the surrounding network, whereas monument noise, local site effects and station-specific artefacts remain in the residual. Whether the prediction or the residual constitutes the useful product depends on the spatial structure of the noise relative to that of the signal, and this dependence is the reason why the evaluation of such methods requires explicit tests of signal preservation. The same principle underlies the removal of tilt and compliance noise from ocean-bottom seismometer data using horizontal and pressure channels \citep{crawfordwebb2000,janiszewski2019} and the self-supervised denoising of distributed acoustic sensing records \citep{batson2019,vandenende2023}, which suggests that the formulation is well founded but that its application to dense geodetic networks has not been examined systematically.

The objectives of the proposed study are threefold. The first is to construct an analysis-ready, leakage-controlled data set of daily displacements for the California network together with fixed temporal and spatial splits. The second is to define a self-supervised reconstruction benchmark and to evaluate a hierarchy of linear and nonlinear models on it, with signal preservation as a primary criterion. The third is to establish whether a deep network provides a measurable improvement over linear spatial filtering and, if it does, to characterise the noise and signal regimes in which the improvement arises. The resulting reconstruction model and its calibrated residuals are intended to serve as the basis for a subsequent continual, near-real-time transient detector.

\section{Literature Review}

\subsection{Common-mode error and network filtering}

Regional filtering of GNSS networks was introduced for the Southern California Integrated GPS Network, where the common-mode error was estimated as the weighted mean of the residuals of the other stations and subtracted from each series \citep{wdowinski1997}. Principal component analysis and the Karhunen--Lo\`eve expansion generalised this approach to spatially varying common modes \citep{dong2006}, and variational Bayesian independent component analysis provided a framework for blind source separation that accommodates missing data \citep{gualandi2016}. Robust trend estimation \citep{blewitt2016} and greedy automatic trajectory decomposition \citep{bedfordbevis2018} address the related problems of velocity estimation and the separation of steps, seasonal terms and transients at individual stations. The current operational standard at NGL is the common-mode component imaging method of \citet{kreemer2021}, which estimates a spatially varying common mode for each station from robust statistics of its neighbours and reduced residual scatter by 39--50\% across several thousand stations in western Europe. Statistical treatments of the offset problem have been compared in a blind community exercise \citep{gazeaux2013}, and the detection of geodetic transients was assessed in the SCEC validation exercise \citep{lohmanmurray2013}. Neither exercise has been repeated with machine-learning methods.

\subsection{Deep learning for geodetic denoising}

Deep-learning approaches to GNSS denoising are recent and few. \citet{mastella2025} trained a single-station model that combines a step detector with a bidirectional recurrent network, using synthetic series whose noise was generated by an adversarial network trained on approximately 5,000 NGL stations; the reported improvement over trajectory-model residuals is established mainly on synthetic data. Network-aware approaches have been developed for Cascadia. \citet{costantino2024} proposed a spatiotemporal graph neural network for the extraction of slow slip signals, evaluated by the recovery of synthetic events and by agreement with tremor activity, which was subsequently used to invert fifteen years of denoised data for the slip history of the subduction interface \citep{costantino2026}. \citet{bachelot2025} compared a graph neural network with stack filtering on two Cascadia data sets and reported common-mode reductions exceeding 70\% and 30\% in the horizontal components, while noting that stack filtering may still be preferable depending on the application. Most recently, \citet{teutschmann2026} pretrained a self-supervised foundation model on 17,652 NGL stations; the model treats stations independently, and its authors identify graph-based spatial encoders, denoising and imputation as directions for future work. To the best of our knowledge, no graph-based or masked-reconstruction model has been trained on the California network, and no study has compared a deep-learning denoiser with the robust common-mode imaging method used operationally by the data provider.

\subsection{Transient detection with synthetic training data}

Because labelled geodetic transients are scarce, the most successful detectors are trained on synthetic signals. \citet{costantino2023} injected 60,000 synthetic slow slip events, computed with elastic dislocation models \citep{okada1992}, into surrogate noise derived from real Cascadia time series and trained a convolutional and attention-based detector on 135 stations. The detector recovered 35 of 40 events in an independent catalogue \citep{michel2019} and reached a detection threshold of approximately $M_w$~6.2--6.4. The authors note that regions with frequent seismicity and overlapping postseismic relaxation would make the problem more difficult and that the approach is region-specific. Alternative strategies include training on real labelled windows, which achieved approximately 75\% accuracy in southwest Japan \citep{tanaka2025}, and convolutional autoencoders for InSAR time series \citep{rouetleduc2021}. In California, the only systematic statewide analysis of transient deformation remains the classical study of \citet{klein2019}, which used approximately 1,000 stations over 1999--2018 and distinguished postseismic, hydrological, anthropogenic and magmatic signals. A recurring weakness of the detection literature is the absence of false-alarm rates measured on real data.

\subsection{Precursory deformation and forecasting}

The geodetic evidence for precursory deformation is contested. The hour-scale acceleration reported by \citet{bletery2023} from 3,026 high-rate series before 90 earthquakes of $M_w \geq 7$ vanishes after common-mode correction \citep{bradley2023}, a point the original authors acknowledge while arguing that the correction removes tectonic signal \citep{bletery2025}. An independent examination of stacked tiltmeter records before the 2011 Tohoku-oki earthquake found no precursory signal above a detection limit of approximately $M_w$~6.4 \citep{hirose2024}. Months-long, thousand-kilometre-scale motions reported before great subduction earthquakes \citep{bedford2020} have likewise been questioned. The best-supported cases are local and slow, such as the eight-month slow slip event preceding the 2014 Iquique earthquake \citep{socquet2017}, and none has been shown to be predictive in the sense of exceeding the rate at which similar transients occur without a subsequent large earthquake.

In earthquake forecasting, geodetic information has improved skill mainly through static strain-rate maps combined with smoothed seismicity, as in prospective global tests \citep{strader2018,bayona2021}. Time-dependent geodetic inputs have not been tested prospectively against the epidemic-type aftershock sequence (ETAS) model, and relative afterslip moment does not correlate with aftershock productivity across 41 earthquakes \citep{churchill2022}. Neural forecasting models trained on catalogues alone have not outperformed ETAS in California \citep{stockman2024}, and simple models frequently match complex networks \citep{mignan2019,mignan2020}. Data-driven forecasting of postseismic deformation has been demonstrated for individual subduction earthquakes \citep{yamaga2019,kano2020}, but not across multiple events or in a strike-slip setting. Likelihood-based evaluation tools for forecasts are available through pyCSEP \citep{savran2022}.

\subsection{Synthesis}

The literature therefore exhibits three gaps that the proposed study addresses directly. Network-aware deep-learning denoising has been developed and evaluated only in subduction settings. Denoising methods are compared predominantly by scatter reduction; the exceptions, which assess the recovery of synthetic signals \citep{mastella2025,costantino2024} or agreement with tremor \citep{costantino2024}, do not examine how recovery depends on the spatial scale of the transient relative to that of the filter. No study has determined whether nonlinear models improve upon linear spatial prediction when both are trained and tested under identical, leakage-controlled conditions. These gaps are prerequisites for the more ambitious questions of transient detection, time-dependent forecasting and precursor testing, since each of these depends on knowing which signals survive the noise-removal step.

\section{Proposed Method}

\subsection{Data}

The study uses NGL daily IGS20 solutions in tenv3 format for all stations within the region $32.4^\circ$--$42.1^\circ$N and $124.6^\circ$--$114.0^\circ$W with at least 1,095 daily solutions, which yields 1,381 stations covering California and adjacent parts of Nevada, Oregon and Baja California \citep{ngl2026}. The data set comprises $8.16 \times 10^6$ station-days. Approximately 875 stations were operating simultaneously in 2008 and 1,040--1,060 between 2015 and 2022. The median completeness within each station's operating period is 95.5\%. In 2018 the median distance to the nearest station was 9.5~km and the median number of stations within 50~km and 100~km was 23 and 76, respectively. The NGL step file lists 3,564 equipment changes and 5,178 potential seismic offsets for these stations. The analysis is restricted to 2008 onwards, when network density is approximately stationary.

\subsection{Preprocessing and data partitioning}

For each station $i$ and component $c \in \{E, N, U\}$, a trajectory model comprising a constant, a linear rate, annual and semi-annual harmonics and Heaviside steps at catalogued epochs is fitted by least squares,
\begin{equation}
x_{i,c}(t) = a + v\,(t - t_0) + \sum_{k=1}^{2}\left[s_k \sin(2\pi k t) + c_k \cos(2\pi k t)\right] + \sum_{j} b_j\, H(t - t_j) + r_{i,c}(t).
\end{equation}
The model is fitted exclusively on the training period, so that no information from the validation or test periods enters the inputs. Equipment changes are removed using the estimated step amplitudes, whereas seismic offsets are retained because they are spatially coherent and should be predictable from the network. Observations deviating by more than five robust standard deviations from a running median are treated as missing. The residuals $r_{i,c}(t)$, expressed in millimetres together with the formal uncertainties $\sigma_{i,c}(t)$ and a binary availability mask, form a three-dimensional array of stations by days by components. For each station and each exclusion radius $R$ (Section~\ref{sec:exclusion}), the sixteen nearest stations at distances $R < d \leq 300$~km define its neighbourhood graph.

The temporal partition assigns 2008--2019 to training, 2020--2021 to validation and 2022--2026 to testing; the test period is evaluated once. Ten per cent of the stations, chosen to be spatially dispersed, are withheld entirely to assess generalisation to unseen sites. The interval from July 2019 to June 2020 is excluded from training for all stations within 150~km of the Ridgecrest rupture, so that the preservation of its postseismic signal can be tested on data unseen by every model.

\subsection{Self-supervised reconstruction task}

Let $\mathcal{M}$ denote a set of withheld station-day-component entries. A model $f_\theta$ receives the array with the entries in $\mathcal{M}$ removed and returns predictions $\hat{x}_{i,c}(t)$ for them. Models with trainable parameters are fitted by minimising the Gaussian negative log-likelihood over the withheld entries,
\begin{equation}
\mathcal{L}(\theta) = \sum_{(i,t,c) \in \mathcal{M}} \left[ \frac{\left(r_{i,c}(t) - \hat{x}_{i,c}(t)\right)^2}{2\,\sigma_{i,c}^2(t)} + \log \sigma_{i,c}(t) \right],
\end{equation}
so that the formal uncertainties down-weight unreliable observations. Three masking patterns are used: randomly scattered entries comprising 10\% of the data, contiguous 30-day blocks at individual stations, and complete withholding of the held-out stations. The masks applied to the test period are generated once with a fixed seed and shared by all models. The reconstruction $\hat{x}$ is interpreted as the network-coherent component, and the residual $r - \hat{x}$ as the station-specific component. Cleaned series and all signal-preservation scores are computed in leave-one-out fashion, with the target station withheld for the entire period, since a flexible model given access to the target could reproduce it and yield a trivially small residual.

\subsection{Exclusion radius and temporal context}
\label{sec:exclusion}

Leave-one-out prediction removes whatever is coherent at the scale of the station spacing, and real transients such as postseismic relaxation, slow slip and hydrological loading are themselves spatially coherent. The common-mode error, by contrast, is coherent over hundreds of kilometres \citep{kreemer2021}, whereas a transient with footprint $L$ affects stations only to distances of the order of $L$. Neighbours closer than an exclusion radius $R$ are therefore withheld from the prediction, and $R$ is treated as an experimental variable rather than a tuning parameter, with $R \in \{0, 10, 25, 50, 100\}$~km. Every model is fitted and evaluated separately at each radius. Larger radii are expected to preserve more transient signal at the cost of weaker noise suppression, and they also emulate sparser networks: at $R = 50$--$100$~km the California network approaches the station spacing of the publicly available network in T\"urkiye, where the median distance to the nearest station is approximately 60~km.

A second route by which signal can enter the prediction is the target's own history. If a model observes the target on days adjacent to the withheld interval, it can extrapolate the onset of a transient into that interval. Three temporal contexts are therefore distinguished: same-day prediction, which uses neighbours on the same day only; causal prediction, which uses data up to the prediction day; and two-sided prediction, which uses a window extending on both sides. Use of the target's own unwithheld data is a separate factor. It is disabled for cleaned series and signal-preservation tests and is permitted only for the gap-filling scores, in which case results report it explicitly. Models M1--M3 are same-day by construction, M4 is two-sided, and M5 is evaluated in all three contexts with and without own history. The causal context is the one required by the continual detector of Section~\ref{sec:continual}.

\subsection{Model hierarchy}

Six models of increasing capacity, eight variants in total, are compared, summarised in Table~\ref{tab:models}. Model M0 predicts zero and provides the floor. Model M1 is the distance-weighted mean of neighbouring residuals in the manner of regional stacking \citep{wdowinski1997}. Model M2 replaces the mean by a robust, outlier-resistant neighbour statistic in the spirit of common-mode component imaging \citep{kreemer2021}. Model M3 is a ridge regression on the sixteen neighbours, trained on the training period; missing neighbour values are set to zero and accompanied by availability indicators, and neighbours are randomly withheld during fitting, so that a single model accommodates arbitrary gap patterns. Because weights estimated for one station cannot be applied to a station absent from training, M3 is implemented in two variants: M3a estimates a separate weight vector for each station and is evaluated on stations seen during training, whereas M3b estimates a single weight vector shared by all stations, with neighbour residuals interacted with inter-station distance and azimuth, and is therefore applicable to the held-out stations. Model M4 extends M3 to frequency-dependent weights by regressing on band-pass filtered neighbour series, analogous to a multichannel Wiener filter, with the same per-station (M4a) and pooled (M4b) variants.

Model M5 is a graph-attention network. For a target station and its neighbours, a 64-day window of residuals, uncertainties and availability flags is encoded by a temporal convolutional network with weights shared across stations. The target embedding attends to the neighbour embeddings, with attention conditioned on inter-station distance and azimuth, and a decoder predicts the withheld values. The network contains approximately $5 \times 10^5$ parameters. This capacity is chosen deliberately: the common-mode process is a single realisation per day, so the effective number of independent training examples for the network-wide component is of the order of the number of days in the training period, approximately 4,400 for 2008--2019, rather than the number of station-days. Training samples random target stations, windows and masks from the training period, and early stopping is based on the validation period. A separate network is trained for each combination of exclusion radius, temporal context and own-history setting that is evaluated. Training is feasible on a single consumer graphics processor with 8~GB of memory.

\begin{table}[ht]
\centering
\caption{Model hierarchy for the reconstruction benchmark.}
\label{tab:models}
\small
\begin{tabular}{l >{\raggedright\arraybackslash}p{5.2cm} l >{\raggedright\arraybackslash}p{4.6cm}}
\toprule
Model & Description & Trainable & Reference approach \\
\midrule
M0 & Zero prediction & No & Floor \\
M1 & Distance-weighted neighbour mean & No & \citet{wdowinski1997} \\
M2 & Robust neighbour statistic & No & \citet{kreemer2021} \\
M3a & Ridge regression on neighbours, per station & Yes & Linear spatial prediction \\
M3b & Ridge regression on neighbours, pooled over stations with geometric features & Yes & Linear spatial prediction \\
M4a & Band-split ridge regression, per station & Yes & Multichannel Wiener filter \\
M4b & Band-split ridge regression, pooled over stations & Yes & Multichannel Wiener filter \\
M5 & Graph-attention temporal network & Yes & \citet{bachelot2025,costantino2024} \\
\bottomrule
\end{tabular}
\end{table}

\subsection{Evaluation}

Reconstruction accuracy is measured by the root-mean-square error of the withheld entries, normalised by each station's noise level and reported separately for each masking pattern, component and hold-out type. Because scatter reduction alone rewards the removal of genuine signal, two additional criteria are adopted. First, synthetic transients are injected into the test-period data before reconstruction. These comprise spatially smooth displacement fields with Gaussian footprints of 10--100~km, amplitudes of 2--10~mm and durations of 10--90~days, followed in a second phase by elastic dislocation sources on strike-slip and creeping faults \citep{okada1992}. The recovery ratio
\begin{equation}
\rho = \frac{\langle r - \hat{x},\, s \rangle}{\langle s,\, s \rangle},
\end{equation}
where $s$ is the injected signal, quantifies the fraction of the transient retained in the residual, and is mapped as a function of footprint, amplitude, duration and exclusion radius. For each model, the common-mode reduction and the recovery ratio are then plotted against one another as $R$ varies, which summarises the trade-off between noise suppression and signal leakage in a single curve. Second, the fraction of the observed Ridgecrest postseismic displacement retained in the residual is computed for the excluded interval. Velocity repeatability between disjoint time windows provides an additional check on whether long-period signal is distorted. Differences between models are assessed with paired bootstrap resampling over days and stations, and all results are reported with their dispersion across random seeds for the trainable models.

\subsection{Extension to continual transient detection}
\label{sec:continual}

If a reconstruction model is shown to preserve injected transients, its residuals provide the input to a transient detector. A multi-scale cumulative-sum statistic on the residuals, combined with a spatial-coherence criterion distinguishing isolated station anomalies from clusters of neighbouring stations, will be calibrated to a prescribed number of false alarms per station-year on quiet intervals. The detector will then be operated in a strictly chronological replay of 2010--2026 in which the reconstruction model is updated as new data arrive. To prevent slow transients from being absorbed as normal behaviour, model updates will be restricted to windows not flagged as anomalous and regularised against catastrophic forgetting \citep{kirkpatrick2017}. Detection delay as a function of transient size and duration, and the extent to which adaptation suppresses slow signals, will be quantified with injected events. This extension is contingent on the outcome of the benchmark and is described here to indicate the intended trajectory of the research.

\section{Expected Findings}

The study is designed so that each plausible outcome is informative. Three hypotheses are tested.

The first hypothesis is that linear spatial prediction captures most of the common-mode noise. We expect model M3 to reduce the normalised reconstruction error substantially relative to models M1 and M2, because station-specific neighbour weights accommodate the spatial variability of the common mode that a uniform stack cannot represent. If M2 performs comparably to M3, the operational common-mode imaging approach would be confirmed as near-optimal for this network, which would itself be a useful result.

The second hypothesis is that a nonlinear model improves reconstruction only where the relation between stations depends on frequency, on station availability or on transient regimes such as postseismic intervals. We anticipate that M5 will outperform M3 and M4 by a modest margin in reconstruction error, with the largest gains at stations with intermittent neighbours and at periods near the draconitic and fortnightly frequencies of GNSS-specific noise. Whether this improvement is accompanied by equal or better transient recovery is the decisive question. An improvement in reconstruction error that coincides with lower recovery ratios would indicate that the network is over-smoothing, absorbing real signal into the coherent component, and would argue against its use for transient detection.

The third hypothesis is that the recovery of transients depends primarily on the ratio of the transient footprint to the exclusion radius. We expect a transient to be largely absorbed when its footprint $L$ exceeds $R$, so that it reaches the stations used for prediction, and to be largely preserved when $R$ exceeds $L$, with recovery rising from near zero towards unity across the transition. Noise suppression is expected to decline more slowly with $R$, because the common mode is coherent over larger distances than the transients considered, so that an intermediate radius should preserve transients of a given scale while retaining most of the noise reduction. We further expect own history to lower the recovery of transients with short onsets in the deep model, and the comparison with and without it will quantify this temporal leakage. The resulting maps will indicate the spatial scale below which daily GNSS can resolve aseismic transients after common-mode removal, and will thereby provide detection thresholds for subsequent searches for slow deformation in California.

The principal deliverables are an open, leakage-controlled benchmark for network reconstruction of daily GNSS data, including the data array, splits, masks and injected signals; a quantitative comparison of linear and deep models under identical conditions; and maps of transient recovery as a function of spatial and temporal scale. A negative result for the deep model, in which it fails to surpass linear prediction on transient recovery, would be reported with equal weight, since it would establish that the additional complexity is not warranted for this application. In either case, the calibrated residuals form the foundation for the continual detection framework outlined above and, subsequently, for testing time-dependent geodetic covariates in likelihood-based earthquake forecast evaluation.

\bibliographystyle{plainnat}
\bibliography{references}

\end{document}
